\section{Funciones de Activación}
\label{sec:annex-funciones-activacion}

Este anexo documenta la librería completa de funciones de activación
disponibles en el modelo. La librería proporciona \textbf{43 funciones de
activación feed-forward}: 40 activaciones elementales y 3 variantes de FFN con
compuerta. Las activaciones se agrupan en cinco familias, siguiendo la
taxonomía de Dubey et al.~\citep{dubey2021activations} y la visión analítica de
Lederer~\citep{lederer2021activations}, y se añade la \emph{Rational Activation
Function} (RAF) aprendible de Fang et al.~\citep{fang2023raft}.

La selección se controla mediante el campo de configuración
\texttt{ffn\_activation} (un enumerado de cadenas) y se despacha en tiempo de
ejecución mediante una factoría que replica el patrón de las capas de
normalización (Anexo~7).

\[
\texttt{ffn\_activation} \in \{\texttt{silu},\;\texttt{gelu},\;\dots,\;
\texttt{raf},\;\texttt{swiglu},\;\texttt{geglu},\;\texttt{reglu}\}
\]

\subsection{Familia Clásica / Sigmoid--Tanh (12 activaciones)}

La familia clásica comprende activaciones suaves, sin estado, basadas en
funciones logísticas y algebraicas~\citep{lederer2021activations,dubey2021activations}.
Todas comparten monotonicía o casi-monotonicía y son continuamente derivables.

\begin{align}
  \text{Sigmoid}(x) &= \frac{1}{1+e^{-x}}
    & \text{Rango: } (0,1),\;\text{monótona, suave}
    & \quad \citep{lederer2021activations} \\[4pt]
  \text{Tanh}(x) &= \frac{e^{x}-e^{-x}}{e^{x}+e^{-x}} = 2\,\sigma(2x)-1
    & \text{Rango: } (-1,1),\;\text{monótona, suave}
    & \quad \citep{lederer2021activations} \\[4pt]
  \text{Arctan}(x) &= \arctan(x)
    & \text{Rango: } (-\tfrac{\pi}{2},\tfrac{\pi}{2}),\;\text{monótona, suave}
    & \quad \citep{lederer2021activations} \\[4pt]
  \text{Softsign}(x) &= \frac{x}{1+|x|}
    & \text{Rango: } (-1,1),\;\text{suave, deriv.\ una vez en 0}
    & \quad \citep{lederer2021activations} \\[4pt]
  \text{Elliott}(x) &= \frac{x}{1+|x|}
    & \text{Alias de Softsign (notación ``elliottsig'' del survey)}
    & \quad \citep{lederer2021activations} \\[4pt]
  \text{Identity}(x) &= x
    & \text{Rango: } (-\infty,\infty),\;\text{lineal, }C^\infty
    & \quad \text{(interna)}
\end{align}

\begin{align}
  \text{Softplus}(x) &= \log(1+e^{x})
    & \text{Rango: } [0,\infty),\;\text{suave, }C^\infty
    & \quad \citep{glorot2011softplus} \\[4pt]
  \text{Mish}(x) &= x\,\tanh\!\big(\text{softplus}(x)\big)
    = x\,\tanh\!\big(\log(1+e^{x})\big)
    & \text{Rango: } [\approx{-0.31},\infty),\;\text{no monótona, suave}
    & \quad \citep{misra2019mish} \\[4pt]
  \text{GELU}(x) &= x\,\Phi(x) = x\cdot\tfrac{1}{2}\Big(1+\text{erf}\!\Big(\frac{x}{\sqrt{2}}\Big)\Big)
    & \text{Rango: } (\approx{-0.17},\infty),\;\text{suave, probabilística}
    & \quad \citep{hendrycks2016gelu} \\[4pt]
  \text{GELU-tanh}(x) &= \tfrac{x}{2}\Big(1+\tanh\!\Big(\sqrt{\tfrac{2}{\pi}}
    (x+0.044715\,x^{3})\Big)\Big)
    & \text{Aproximación con tanh de GELU (más rápida)}
    & \quad \citep{hendrycks2016gelu} \\[4pt]
  \text{ReLU}(x) &= \max(0,x)
    & \text{Rango: } [0,\infty),\;\text{tramos lineales, no suave en 0}
    & \quad \citep{nair2010relu} \\[4pt]
  \text{SiLU}(x) &= x\,\sigma(x) = \frac{x}{1+e^{-x}}
    & \text{Rango: } (\approx{-0.278},\infty),\;\text{no monótona, suave}
    & \quad \citep{elfwing2018silu}
\end{align}

\noindent\textbf{SiLU} (Sigmoid Linear Unit) es la activación por defecto
(\texttt{ffn\_activation=silu}). Cuando $\beta=1$, el Swish paramétrico se
reduce a SiLU. GELU y su aproximación con tanh son el estándar en BERT y
GPT-2/3~\citep{hendrycks2016gelu}. Mish~\citep{misra2019mish} es una
composición auto-regularizada, no monótona, de Softplus y Tanh, que supera a
Swish en varios benchmarks.

\subsection{Familia Rectificada (12 activaciones)}

La familia de rectificadores extiende ReLU con pendientes aprendibles,
variantes acotadas y formas hiperbólicas
rectificadas~\citep{nair2010relu,dubey2021activations}.

\begin{align}
  \text{LeakyReLU}(x) &= \begin{cases} x & x\ge 0 \\ a\,x & x < 0 \end{cases},
    \quad a=0.01
    & \text{Rango: } (-\infty,\infty)
    & \quad \citep{maas2013leakyrelu} \\[4pt]
  \text{ReLU6}(x) &= \min\!\big(\max(0,x),\,6\big)
    & \text{Rango: } [0,6],\;\text{acotada, amigable para coma fija}
    & \quad \citep{howard2017mobile} \\[4pt]
  \text{HardSwish}(x) &= x\cdot\frac{\text{ReLU6}(x+3)}{6}
    & \text{Rango: } (\approx{-1.67},\infty),\;\text{aprox.\ barata de Swish}
    & \quad \citep{howard2019hardswish} \\[4pt]
  \text{PReLU}(x) &= \max(0,x)+\mathbf{p}\odot\min(0,x)
    & \text{Pendiente por canal aprendible } \mathbf{p},\;\text{init 0.25}
    & \quad \citep{he2015prelu} \\[4pt]
  \text{AbsReLU}(x) &= \max(0,x)-\max(0,-x) = x\;\text{rectificado a rampa de signo}
    & \text{Rango: } (-\infty,\infty)
    & \quad \citep{dubey2021activations} \\[4pt]
  \text{NLReLU}(x) &= \beta\,\log\!\big(1+\max(0,x)\big)
    & \text{Rango: } [0,\infty),\;\beta=1.0,\;\text{compresión logarítmica}
    & \quad \citep{dubey2021activations}
\end{align}

\begin{align}
  \text{BReLU}(x) &= \min\!\big(\max(0,x),\,t\big),\quad t=1.0
    & \text{Rango: } [0,t],\;\text{rectificador acotado}
    & \quad \citep{dubey2021activations} \\[4pt]
  \text{VReLU}(x) &= |x|
    & \text{Rango: } [0,\infty),\;\text{rectificador simétrico}
    & \quad \citep{dubey2021activations} \\[4pt]
  \text{Hexpo}(x) &= a\,\max(0,x)-c\,\max(0,-x),\quad a{=}c{=}1.0
    & \text{Rectificador asimétrico bilateral}
    & \quad \citep{dubey2021activations} \\[4pt]
  \text{PenalizedTanh}(x) &= \max\!\big(0,\,\tanh(x)\big)
    & \text{Rango: } [0,1),\;\text{tanh rectificado, suave cerca de 0}
    & \quad \citep{dubey2021activations} \\[4pt]
  \text{DisReLU}(x) &= \max(0,\,x-\delta),\quad \delta=0.0
    & \text{ReLU desplazado a la derecha}
    & \quad \citep{dubey2021activations} \\[4pt]
  \text{LiSHT}(x) &= x\,\tanh(x)
    & \text{Rango: } [0,\infty)\;\text{(magnitud), no monótona, suave}
    & \quad \citep{roy2019lisht}
\end{align}

\noindent\textbf{PReLU}~\citep{he2015prelu} convierte la pendiente negativa en
un parámetro aprendible por canal, permitiendo a la red aprender la fuga
óptima por característica. \textbf{HardSwish}~\citep{howard2019hardswish} es
la aproximación barata de Swish usada en MobileNetV3.
\textbf{ReLU6}~\citep{howard2017mobile} se introdujo en MobileNetV1 por su
compatibilidad con cuantización de coma fija.

\subsection{Familia Exponencial / ELU (12 activaciones)}

La familia exponencial-lineal proporciona saturación negativa suave y
variantes auto-normalizadoras~\citep{clevert2016elu,klambauer2017selu,dubey2021activations}.
Varios miembros introducen parámetros aprendibles por canal.

\begin{align}
  \text{ELU}(x) &= \begin{cases} x & x\ge 0 \\ \alpha(e^{x}-1) & x<0 \end{cases},
    \quad \alpha=1.0
    & \text{Rango: } (-\alpha,\infty)
    & \quad \citep{clevert2016elu} \\[4pt]
  \text{SELU}(x) &= \lambda\,\text{ELU}_{\alpha'}(x)
    & \alpha'{\approx}1.6733,\;\lambda{\approx}1.0507,
    & \quad \citep{klambauer2017selu} \\
  & & \text{auto-normalizadora (media 0, var 1)} & \notag \\[4pt]
  \text{CELU}(x) &= \begin{cases} x & x\ge 0 \\ \alpha(e^{x/\alpha}-1) & x<0 \end{cases},
    \quad \alpha=1.0
    & \text{Derivable una vez en 0 para cualquier } \alpha
    & \quad \citep{barron2017celu} \\[4pt]
  \text{PELU}(x) &= \begin{cases} \tfrac{\alpha}{\beta}\,x & x\ge 0 \\
    \alpha(e^{x/\beta}-1) & x<0 \end{cases}
    & \alpha,\beta\text{ aprendibles por canal}
    & \quad \citep{trottier2017pelu} \\[4pt]
  \text{MPELU}(x) &= \text{ReLU}(x)+\beta\cdot\text{ELU}_{\alpha}(x)\cdot\mathbb{1}_{x<0}
    & \alpha,\beta\text{ aprendibles por canal}
    & \quad \citep{dubey2021activations}
\end{align}

\begin{align}
  \text{FELU}(x) &= \begin{cases} x & x\ge 0 \\ \alpha(e^{x}-1) & x<0 \end{cases}
    & \alpha\text{ aprendible, acotado a } [0,1]
    & \quad \citep{dubey2021activations} \\[4pt]
  \text{EELU}(x) &= \beta\,\text{ReLU}(x)+\alpha(e^{x}-1)\cdot\mathbb{1}_{x<0}
    & \alpha,\beta\text{ aprendibles por canal}
    & \quad \citep{dubey2021activations} \\[4pt]
  \text{PDELU}(x) &= \begin{cases} x & x\ge 0 \\
    \alpha(e^{x/\alpha}-1)+(1-\alpha)\,x & x<0 \end{cases}
    & \alpha\text{ aprendible, interpola CELU--identidad}
    & \quad \citep{dubey2021activations} \\[4pt]
  \text{PREU}(x) &= \beta\,\text{ReLU}(x)+\alpha(e^{x}-1)\cdot\mathbb{1}_{x<0}^{\le 0}
    & \alpha,\beta\text{ aprendibles por canal}
    & \quad \citep{dubey2021activations} \\[4pt]
  \text{SoftExp}(x) &= \begin{cases}
    \tfrac{1}{\alpha}\big(e^{\alpha x}-1\big)+\alpha & \alpha>0 \\
    x & \alpha=0 \\
    -\tfrac{1}{\alpha}\ln\!\big(1-\alpha(x+\alpha)\big) & \alpha<0
    \end{cases}
    & \alpha\text{ aprendible; interp.\ exp/lineal/log}
    & \quad \citep{godfrey2015softexp}
\end{align}

\begin{align}
  \text{ELiSH}(x) &= \begin{cases} x\,\sigma(x) & x\ge 0 \\
    (e^{x}-1)\,\sigma(x) & x<0 \end{cases}
    & \text{Rama pos.\ Swish, rama neg.\ con compuerta ELU}
    & \quad \citep{basirat2019elish} \\[4pt]
  \text{HardELiSH}(x) &= \begin{cases} x\cdot\text{hard\_sigmoid}(x) & x\ge 0 \\
    (e^{x}-1)\cdot\text{hard\_sigmoid}(x) & x<0 \end{cases}
    & \text{Variante con sigmoide rígida de ELiSH}
    & \quad \citep{basirat2019elish}
\end{align}

\noindent donde $\text{hard\_sigmoid}(x)=\text{ReLU6}(x+3)/6$.

\noindent\textbf{ELU}~\citep{clevert2016elu} aporta saturación negativa suave.
\textbf{SELU}~\citep{klambauer2017selu} induce auto-normalización con
$\alpha'{\approx}1.6733$, $\lambda{\approx}1.0507$ fijos.
\textbf{CELU}~\citep{barron2017celu} es derivable una vez en 0 para cualquier
$\alpha$ (a diferencia de ELU). \textbf{PELU}~\citep{trottier2017pelu} hace
ambos $\alpha,\beta$ aprendibles por canal. Las variantes paramétricas
(MPELU, FELU, EELU, PDELU, PREU) generalizan ELU y PReLU introduciendo
parámetros aprendibles por canal para las ramas negativa y positiva.
\textbf{SoftExp}~\citep{godfrey2015softexp} interpola continuamente entre los
regímenes logarítmico ($\alpha<0$), lineal ($\alpha=0$) y exponencial
($\alpha>0$). \textbf{ELiSH} y \textbf{HardELiSH}~\citep{basirat2019elish}
combinan comportamientos de Swish y ELU mediante una definición a tramos.

\subsection{Swish y Activaciones Paramétricas}

Swish~\citep{ramachandran2017swish} es una familia parametrizada por una
pendiente $\beta$:
\[
  \text{Swish}_\beta(x) = x\,\sigma(\beta x) = \frac{x}{1+e^{-\beta x}}.
\]
Cuando $\beta=1$ se obtiene SiLU (el valor por defecto). Con $\beta$ grande se
recupera ReLU; con $\beta\to 0$ se aproxima a la función lineal $x/2$.

\begin{itemize}[leftmargin=*]
  \item \textbf{Swish} (\texttt{swish}): $\beta$ fijo (por defecto 1.0,
    configurado mediante \texttt{ffn\_activation\_config.swish\_beta}). Sin
    parámetros aprendibles.
  \item \textbf{SwishTrainable} (\texttt{swish\_trainable}): $\beta$ es un
    parámetro aprendible por canal, inicializado a 1.0. Equivalente a SiLU al
    inicio, se adapta durante el entrenamiento.
  \item \textbf{Maxout} (\texttt{maxout}): calcula $k$ proyecciones lineales y
    toma el máximo elemento a elemento: $\max_{i=0}^{k-1}(W_i x + b_i)$.
    Aproximador universal de cualquier función continua con suficientes piezas.
    El número de parámetros aumenta por un factor $k$ (por defecto $k=2$,
    configurado mediante \texttt{ffn\_activation\_config.maxout\_pieces}).
\end{itemize}

\noindent Referencias: Swish~\citep{ramachandran2017swish},
Maxout~\citep{goodfellow2013maxout}.

\subsection{Funciones de Activación Racionales (RAF)}
\label{sec:raf-es}

Siguiendo a Fang et al.~\citep{fang2023raft}, una \emph{Rational Activation
Function} es una razón aprendible de dos polinomios de bajo grado (un
aproximante de Pad\'e):
\[
  F(x) = \frac{P(x)}{Q(x)}
        = \frac{\displaystyle\sum_{j=0}^{m} a_j\, x^{j}}
               {1 + R(x)}
\]
donde $\mathbf{a}\in\mathbb{R}^{m+1}$ y $\mathbf{b}\in\mathbb{R}^{n}$ son
\emph{aprendibles}, y el t\'ermino del denominador $R(x)$ depende de la
versión. La configuración por defecto coincide con el artículo:
grado $(m,n)=(5,4)$, versión ``A'' (segura), inicializada mediante un ajuste
por mínimos cuadrados a GELU en $[-3,3]$.

La implementación admite cinco variantes de denominador:

\begin{itemize}[leftmargin=*]
  \item \textbf{Versión A} (por defecto, ``segura''):
    $Q(x)=1+\sum_{k}|b_k\, x^{k+1}|$ (valor absoluto por t\'ermino).
    Garantiza $Q(x)\ge 1$, sin división por cero.
  \item \textbf{Versión B}: $Q(x)=1+\big|\sum_{k} b_k\, x^{k+1}\big|$
    (valor absoluto de la suma completa). También garantiza $Q(x)\ge 1$.
  \item \textbf{Versión C}: $Q(x)=0.1+\big|\sum_{k} b_k\, x^{k}\big|$
    con un suelo de $0.1$. Denominador mínimo $0.1$.
  \item \textbf{Versión D}: como B, pero inyecta ruido multiplicativo uniforme
    $U(1-\varepsilon, 1+\varepsilon)$ sobre los pesos del denominador
    \emph{sólo durante el entrenamiento} ($\varepsilon=0.1$). Efecto de
    regularización.
  \item \textbf{Versión N} (``no segura''):
    $Q(x)=1+\sum_{k} b_k\, x^{k+1}$ sin valor absoluto.
    Puede tener polos (ceros del denominador); usar con cuidado.
\end{itemize}

\textbf{Escalado de entrada RAFT.} Dado que un aproximante de Pad\'e es un
aproximador universal fiable sólo en un intervalo acotado, el modelo RAFT
preprocesa las pre-activaciones de cada token con un escalado min--max por
token hacia $[-3,3]$:
\[
  \tilde{x} = \frac{x-\min(x)}{\max(x)-\min(x)}\cdot 6 - 3.
\]
Se activa con \texttt{ffn\_activation\_config.raf\_input\_scaling}
(por defecto \texttt{false}).

\textbf{Objetivos de inicialización.} Los coeficientes racionales pueden
inicializarse ajustándose a cualquiera de las siguientes funciones objetivo:
\texttt{gelu} (por defecto), \texttt{relu}, \texttt{leaky\_relu},
\texttt{leaky\_relu\_0.1}, \texttt{sigmoid}, \texttt{tanh},
\texttt{swish}/\texttt{silu}, o \texttt{identity}. El ajuste usa una
resolución por mínimos cuadrados sobre 2001 muestras uniformemente espaciadas
en $[-3,3]$.

\textbf{Congelado.} Fijar \texttt{raf\_trainable=false} congela los parámetros
racionales, útil para ajuste eficiente en parámetros. Las RAF congeladas se
emparejan típicamente con una tasa de aprendizaje separada (mayor) para el
resto de parámetros de la red.

\textbf{Nota sobre el nombre.} El acrónimo RAF significa \emph{Rational}
Activation Function, no ``Rectified''; corresponde a la razón polinómica
aprendible descrita arriba.

\subsection{Variantes de FFN con Compuerta: SwiGLU, GEGLU, ReGLU}
\label{sec:glu-es}

SwiGLU, GEGLU y ReGLU~\citep{shazeer2020glu} son unidades feed-forward con
compuerta que reemplazan el bloque estándar
$\text{Linear}\to\text{activación}\to\text{Linear}$. Poseen sus propias
proyecciones lineales y, por tanto, se implementan como módulos FFN, \emph{no}
como activaciones elementales:
\[
  \text{GatedFFN}(x) = \text{dropout}\!\Big(\,
    W_{\text{down}}\Big(\text{act}(x\,W_{\text{gate}})
    \odot (x\,W_{\text{up}})\Big)\Big),
\]
donde $W_{\text{gate}}, W_{\text{up}}\in\mathbb{R}^{d\times d_{\text{ff}}}$ y
$W_{\text{down}}\in\mathbb{R}^{d_{\text{ff}}\times d}$. La función de compuerta
es:
\begin{itemize}[leftmargin=*]
  \item \textbf{SwiGLU} (\texttt{swiglu}): $\text{act}=\text{SiLU}$
    (por defecto, usado en Llama, PaLM).
  \item \textbf{GEGLU} (\texttt{geglu}): $\text{act}=\text{GELU}$
    (usado en algunas variantes de T5).
  \item \textbf{ReGLU} (\texttt{reglu}): $\text{act}=\text{ReLU}$
    (la opción más barata).
\end{itemize}

\noindent Cuando \texttt{ffn\_activation} es uno de
\texttt{swiglu}/\texttt{geglu}/\texttt{reglu}, la capa feed-forward sustituye
los bloques FFN densos y MoE por un FFN con compuerta construido con la misma
clase de proyección (BitLinear bajo BitNet). El sesgo está desactivado por
defecto.

\subsection{Esquema de Configuración}

El objeto \texttt{ffn\_activation\_config} agrupa todos los parámetros de las
activaciones aprendibles. Es opcional y se ignora para activaciones sin estado.

\begin{itemize}[leftmargin=*]
  \item \texttt{raf\_degrees}: $[m,n]$ grados numerador/denominador del Pad\'e
    (enteros $\ge 1$). Por defecto: $[5,4]$.
  \item \texttt{raf\_version}: forma del denominador, uno de
    \texttt{A} (por defecto), \texttt{B}, \texttt{C}, \texttt{D}, \texttt{N}.
  \item \texttt{raf\_approx\_func}: objetivo de inicialización. Uno de
    \texttt{gelu} (por defecto), \texttt{relu}, \texttt{leaky\_relu},
    \texttt{leaky\_relu\_0.1}, \texttt{sigmoid}, \texttt{tanh},
    \texttt{swish}, \texttt{silu}, \texttt{identity}.
  \item \texttt{raf\_trainable}: congela los parámetros racionales cuando es
    \texttt{false}. Por defecto: \texttt{true}.
  \item \texttt{raf\_input\_scaling}: escalado min--max por token a $[-3,3]$
    antes del racional (preprocesamiento RAFT). Por defecto: \texttt{false}.
  \item \texttt{prelu\_init}: pendiente negativa inicial para PReLU.
    Por defecto: $0.25$.
  \item \texttt{elu\_alpha}: constante de saturación de ELU. Por defecto: $1.0$.
  \item \texttt{celu\_alpha}: constante de saturación de CELU. Por defecto: $1.0$.
  \item \texttt{swish\_beta}: $\beta$ fijo para Swish / $\beta$ inicial para
    SwishTrainable. Por defecto: $1.0$.
  \item \texttt{leaky\_relu\_slope}: pendiente negativa para LeakyReLU.
    Por defecto: $0.01$.
  \item \texttt{maxout\_pieces}: número de piezas lineales $k$ para Maxout
    (entero $\ge 1$). Por defecto: $2$.
  \item Claves adicionales de la familia ELU paramétrica:
    \texttt{pelu\_alpha}, \texttt{mpelu\_alpha}, \texttt{mpelu\_beta},
    \texttt{felu\_alpha}, \texttt{eelu\_alpha}, \texttt{eelu\_beta},
    \texttt{pdelu\_alpha}, \texttt{preu\_alpha}, \texttt{preu\_beta},
    \texttt{softexp\_alpha}. Todas por defecto $1.0$ (excepto
    \texttt{softexp\_alpha} que por defecto es $0.0$).
\end{itemize}

\subsection{Guía de Selección}

El siguiente árbol de decisión ayuda a elegir la activación:
\begin{itemize}[leftmargin=*]
  \item \textbf{Por defecto seguro} $\to$ \texttt{silu} (SiLU/Swish$_1$, usado
    en Llama, PaLM).
  \item \textbf{Estabilidad clásica} $\to$ \texttt{gelu} o \texttt{relu}.
  \item \textbf{No monótona suave} $\to$ \texttt{mish} o \texttt{swish}.
  \item \textbf{Amigable para móvil} $\to$ \texttt{hardswish} o \texttt{relu6}.
  \item \textbf{Forma por tarea aprendible} $\to$ \texttt{raf},
    \texttt{prelu}, \texttt{swish\_trainable}, o \texttt{maxout}.
  \item \textbf{Redes auto-normalizadoras} $\to$ \texttt{selu} (con
    AlphaDropout).
  \item \textbf{FFN con compuerta (proyecciones propias)} $\to$
    \texttt{swiglu}, \texttt{geglu}, o \texttt{reglu}.
\end{itemize}

\noindent Con cuantización BitNet, \texttt{silu} puede ser más robusta que las
activaciones suaves que dependen de un cálculo de gradiente preciso. Las
activaciones aprendibles (\texttt{raf}, \texttt{prelu},
\texttt{swish\_trainable}) añaden parámetros pero pueden adaptarse a la
distribución de la tarea.